\documentclass[a4paper,12pt]{article}

\usepackage[T2A]{fontenc}
\usepackage[utf8]{inputenc}
\usepackage[russian]{babel}

\usepackage{amsmath,amssymb,amsthm,mathtools}
\usepackage{helvet}
\renewcommand{\familydefault}{\sfdefault}
\usepackage{icomma}
\usepackage[a4paper,margin=2.2cm]{geometry}
\usepackage{microtype}
\usepackage{tikz}
\usepackage{tkz-euclide}
\usepackage{pgfplots}
\pgfplotsset{compat=1.18}
\usepackage{tabularx,booktabs,longtable}
\usepackage{enumitem}
\usepackage{hyperref}
\usepackage{parskip}
\usepackage{fancyhdr}
\usepackage{setspace}
\usepackage{xcolor}
\usepackage{array}

\definecolor{accent}{HTML}{008080}
\definecolor{darkaccent}{HTML}{004D4D}
\definecolor{textgray}{HTML}{333333}
\definecolor{lightaccent}{HTML}{EAF6F6}

\hypersetup{
    colorlinks=true,
    linkcolor=darkaccent,
    urlcolor=darkaccent,
    citecolor=darkaccent
}

\pagestyle{fancy}
\fancyhf{}
\fancyhead[L]{\textcolor{darkaccent}{Чек-лист по математике}}
\fancyhead[R]{\textcolor{darkaccent}{ОГЭ}}
\fancyfoot[C]{\thepage}

\setlength{\headheight}{15pt}
\onehalfspacing

\setlist[itemize]{leftmargin=1.8em}
\setlist[enumerate]{leftmargin=2em}

\newcommand{\topic}{Задачи на производительность и базовые геометрические приёмы}
\newcommand{\studentname}{Ксении Клыковой}

\begin{document}

% ---------------- TITUL ----------------
\begin{titlepage}
\thispagestyle{empty}
\begin{center}
\vspace*{2.2cm}

{\Huge\bfseries \textcolor{darkaccent}{Чек-лист}}\\[0.6cm]
{\LARGE\bfseries \textcolor{textgray}{\topic}}\\[1.5cm]

{\Large Лёвин Артём Александрович\\[0.2cm]
эксклюзивно для \studentname}\\[1.2cm]

{\large \today}

\vfill

\begin{tikzpicture}[x=1cm,y=1cm]
\draw[line width=1.3pt, color=accent]
(0,0) .. controls (2,0.9) and (4,-0.6) .. (6,0.3)
      .. controls (8,1.1) and (10,-0.4) .. (12,0.5);
\draw[line width=0.7pt, color=accent!55]
(0.4,-0.2) .. controls (2.2,0.5) and (4.4,-0.8) .. (6.2,0.05)
          .. controls (8.3,0.9) and (10.1,-0.2) .. (11.8,0.2);
\end{tikzpicture}

\vspace{0.4cm}
{\small\textcolor{textgray}{Сдержанный конспект для повторения, закрепления и самопроверки}}
\end{center}
\end{titlepage}

% ---------------- MAIN ----------------
\section*{1. Тема занятия и её ядро}

На занятии отрабатывались два блока:

\begin{enumerate}
    \item \textbf{Текстовые задачи на производительность}  
    (таблица ``производительность -- время -- работа'', перевод условия в уравнение, контроль смысла разности времен).
    \item \textbf{Геометрические задачи ОГЭ}, где фигура часто \textbf{сводится к треугольнику}:  
    параллелограмм, ромб, квадрат в прямоугольном треугольнике, подобие, теорема синусов, площадь через две стороны и угол между ними.
\end{enumerate}

\section*{2. Задачи на производительность}

\subsection*{2.1. Базовая схема}

Используется связь
\[
\text{работа}=\text{производительность}\cdot \text{время},
\qquad
t=\frac{A}{v},
\]
где $A$ --- объём работы, $v$ --- производительность, $t$ --- время.

\subsection*{2.2. Стандартная таблица}

\renewcommand{\arraystretch}{1.35}
\begin{table}[h!]
\centering
\caption{Шаблон для задач на трубы, насосы, станки, рабочих}
\begin{tabularx}{0.9\linewidth}{>{\raggedright\arraybackslash}p{3.1cm}XXX}
\toprule
\textbf{Объект} & \textbf{Производительность} & \textbf{Время} & \textbf{Работа} \\
\midrule
Первый & $v_1$ & $\dfrac{A_1}{v_1}$ & $A_1$ \\
Второй & $v_2$ & $\dfrac{A_2}{v_2}$ & $A_2$ \\
\bottomrule
\end{tabularx}
\end{table}

\subsection*{2.3. Как выбирать неизвестное}

Если в условии сказано:

\begin{itemize}
    \item ``первая труба пропускает на $15$ л/мин меньше, чем вторая'',  
    удобно взять \textbf{производительность второй} за $x$, тогда первой --- $x-15$;
    \item ``вторая на $6$ л/мин больше первой'',  
    удобно взять \textbf{производительность первой} за $x$, тогда второй --- $x+6$.
\end{itemize}

\subsection*{2.4. Самая важная мысль занятия}

\textbf{Нельзя механически вычитать дроби по шаблону.}  
Сначала нужно понять, \textbf{кто работал дольше}, а уже потом писать уравнение.

Если первая труба медленнее и заполняет \textbf{тот же} резервуар, то её время больше:
\[
\frac{100}{x-15}-\frac{100}{x}=6.
\]

Но если первая труба медленнее, \textbf{зато объём у неё меньше}, то ситуация может измениться: именно смысл текста решает, какая дробь больше.

\subsection*{2.5. Алгоритм решения}

\begin{enumerate}
    \item Выберите неизвестное.
    \item Заполните таблицу.
    \item Выразите время через $\dfrac{\text{работа}}{\text{производительность}}$.
    \item Определите по смыслу, кто работал дольше.
    \item Составьте уравнение на разность времен.
    \item Проверьте ограничение: производительность должна быть \textbf{положительной}.
\end{enumerate}

\subsection*{2.6. Мини-пример}

Первая труба пропускает на $15$ л/мин меньше второй. Резервуар $100$ л вторая заполняет на $6$ минут быстрее. Найти производительность второй.

Пусть производительность второй трубы равна $x$ л/мин. Тогда у первой --- $x-15$ л/мин.

Время:
\[
t_1=\frac{100}{x-15}, \qquad t_2=\frac{100}{x}.
\]

Так как первая труба медленнее, она работает дольше:
\[
\frac{100}{x-15}-\frac{100}{x}=6.
\]

\textbf{Ключевая проверка смысла:} знаменатель $x-15$ меньше, значит дробь $\dfrac{100}{x-15}$ больше.

\subsection*{2.7. Типичные ошибки}

\begin{itemize}
    \item Перепутать, от какого времени что вычитать.
    \item Смотреть только на таблицу и не читать смысл фразы ``быстрее/медленнее'', ``больше/меньше времени''.
    \item Забыть, что разные объёмы работы могут поменять сравнение времён.
    \item Получить отрицательную или слишком малую производительность и не заметить это.
\end{itemize}

% ---------------- GEOMETRY ----------------
\section*{3. Геометрия: ключевая идея занятия}

\textbf{Многие задачи про четырёхугольники на самом деле решаются через треугольники.}

На занятии встречались такие приёмы:

\begin{itemize}
    \item свести параллелограмм к треугольнику по диагонали;
    \item использовать \textbf{теорему синусов}, если в треугольнике известны углы и нужно отношение сторон;
    \item использовать \textbf{подобие} прямоугольных треугольников;
    \item использовать \textbf{площадь} разными способами, чтобы найти неизвестную высоту или расстояние.
\end{itemize}

\section*{4. Параллелограмм и теорема синусов}

\subsection*{4.1. Теорема синусов}

В любом треугольнике
\[
\frac{a}{\sin \alpha}=\frac{b}{\sin \beta}=\frac{c}{\sin \gamma}.
\]

Её особенно удобно применять, когда:

\begin{itemize}
    \item известны два угла;
    \item требуется отношение сторон;
    \item стороны и углы ``смешаны'' в одном вопросе.
\end{itemize}

\subsection*{4.2. Сюжет из занятия}

Диагональ параллелограмма делит угол на $30^\circ$ и $45^\circ$. Нужно найти отношение сторон.

После проведения диагонали задача сводится к треугольнику. В нужном треугольнике стороны параллелограмма становятся сторонами треугольника, а углы известны.

Если в треугольнике $ADC$
\[
\angle ACD=30^\circ,\qquad \angle CAD=45^\circ,
\]
то по теореме синусов
\[
\frac{AD}{\sin 30^\circ}=\frac{CD}{\sin 45^\circ}.
\]

Значит,
\[
\frac{AD}{CD}=\frac{\sin 30^\circ}{\sin 45^\circ}
=\frac{\tfrac12}{\tfrac{\sqrt2}{2}}
=\frac{1}{\sqrt2}
=\frac{\sqrt2}{2}.
\]

\subsection*{4.3. Табличные значения, которые нужны постоянно}

\[
\sin 30^\circ=\frac12,\qquad
\sin 45^\circ=\frac{\sqrt2}{2},\qquad
\tg 45^\circ=1.
\]

\subsection*{4.4. Короткий вывод}

Если в задаче \textbf{даны углы, а спрашивают стороны или их отношение}, очень часто стоит проверять \textbf{теорему синусов}.

\begin{center}
\begin{tikzpicture}[scale=1]
\tkzDefPoint(0,0){A}
\tkzDefPoint(4.8,0){D}
\tkzDefPoint(1.4,2.2){B}
\tkzDefPoint(6.2,2.2){C}
\tkzDrawPolygon[thick,color=accent](A,B,C,D)
\tkzDrawSegment[thick,color=darkaccent](A,C)
\tkzLabelPoints[below](A,D)
\tkzLabelPoints[above](B,C)
\tkzMarkAngle[size=0.45](D,A,C)
\tkzMarkAngle[size=0.7](C,A,B)
\tkzLabelAngle[pos=0.65](D,A,C){\small $45^\circ$}
\tkzLabelAngle[pos=1.02](C,A,B){\small $30^\circ$}
\end{tikzpicture}
\end{center}

% ---------------- SIMILARITY ----------------
\section*{5. Квадрат в прямоугольном треугольнике: подобие и тангенс}

\subsection*{5.1. Суть конструкции}

В прямоугольном треугольнике размещён квадрат так, что две его стороны лежат на катетах. Возникают два прямоугольных треугольника, которые оказываются \textbf{подобными}.

\subsection*{5.2. Признак подобия}

Если два угла одного треугольника соответственно равны двум углам другого, то треугольники подобны.

\subsection*{5.3. Практический вывод}

Если сторона квадрата равна $x$, а прилегающие отрезки на катетах равны $a$ и $b$, то из подобия или из равенства тангенсов получается
\[
\frac{x}{b}=\frac{a}{x}.
\]

Отсюда
\[
x^2=ab.
\]

Значит, \textbf{площадь квадрата}
\[
S=x^2=ab.
\]

\subsection*{5.4. Почему это удобно}

Иногда спрашивают именно площадь квадрата, а не его сторону. Тогда можно сразу использовать
\[
S=ab,
\]
не извлекая корень без необходимости.

\subsection*{5.5. Типичная ошибка}

Получив $x^2=ab$, ученики иногда пишут ответ $x=\sqrt{ab}$, хотя в задаче просили \textbf{площадь квадрата}.  
Если требуется площадь, ответом будет именно
\[
ab.
\]

% ---------------- RHOMBUS ----------------
\section*{6. Ромб: расстояние от центра до стороны}

\subsection*{6.1. Что такое расстояние от точки до стороны}

Это длина \textbf{перпендикуляра}, опущенного из точки на прямую, содержащую сторону.

\subsection*{6.2. Полезные свойства ромба}

\begin{itemize}
    \item все стороны равны;
    \item диагонали пересекаются и делят друг друга пополам;
    \item диагонали ромба являются биссектрисами его углов.
\end{itemize}

\subsection*{6.3. Главный приём занятия}

Вместо прямого поиска высоты удобно сравнить \textbf{две формулы площади одного и того же треугольника}.

Если сторона ромба равна $4$, острый угол равен $30^\circ$, а $O$ --- точка пересечения диагоналей, то удобно рассмотреть треугольник $AOB$.

С одной стороны, его площадь равна половине площади треугольника $ABD$:
\[
S_{AOB}=\frac12 S_{ABD}.
\]

А
\[
S_{ABD}=\frac12 \cdot AB \cdot AD \cdot \sin \angle BAD
= \frac12\cdot 4\cdot 4\cdot \sin 30^\circ = 4.
\]

Следовательно,
\[
S_{AOB}=2.
\]

С другой стороны, если $OH$ --- расстояние от центра ромба до стороны $AB$, то
\[
S_{AOB}=\frac12 \cdot AB \cdot OH
=\frac12\cdot 4\cdot OH=2OH.
\]

Отсюда
\[
2OH=2,\qquad OH=1.
\]

\subsection*{6.4. Ещё одна полезная формула}

Для ромба со стороной $a$ и углом $\alpha$ расстояние от центра до стороны равно
\[
r=\frac{a\sin\alpha}{2}.
\]

Проверка для $a=4$, $\alpha=30^\circ$:
\[
r=\frac{4\cdot \tfrac12}{2}=1.
\]

% ---------------- AREAS ----------------
\section*{7. Формулы площади, которые были нужны}

\begin{enumerate}
    \item Через основание и высоту:
    \[
    S=\frac12 ah
    \quad \text{(для треугольника).}
    \]
    \item Через две стороны и угол между ними:
    \[
    S=\frac12 ab\sin\gamma.
    \]
    \item Площадь одной и той же фигуры можно считать \textbf{разными способами} и приравнивать результаты.
\end{enumerate}

\section*{8. Чек-лист самопроверки перед ответом}

\begin{itemize}
    \item Я точно понял(а), \textbf{что взято за $x$}?
    \item Я не перепутал(а), \textbf{кто работает дольше}, а кто быстрее?
    \item В текстовой задаче я проверил(а) смысл уравнения, а не просто ``подставил(а) по шаблону''?
    \item В геометрии я свёл(свела) фигуру к \textbf{подходящему треугольнику}?
    \item Если есть углы и стороны, я подумал(а) о \textbf{теореме синусов}?
    \item Если есть похожие прямоугольные треугольники, я проверил(а) \textbf{подобие}?
    \item Если нужно расстояние или высота, я попробовал(а) \textbf{площадь двумя способами}?
    \item Я отвечаю именно на то, что спросили: \textbf{сторону, площадь, отношение, расстояние}?
\end{itemize}

\section*{9. Типичные ошибки по этому занятию}

\begin{itemize}
    \item Перепутать ``быстрее'' и ``дольше''.
    \item Неверно определить, какая из дробей времени больше.
    \item Писать уравнение, не сравнив смысл двух ситуаций.
    \item В геометрии работать сразу с четырёхугольником и не заметить полезный треугольник.
    \item Не увидеть подобие.
    \item Подставить не тот угол в формулу площади.
    \item Перепутать сторону квадрата и его площадь.
    \item Записать числовой ответ с неверной запятой.
\end{itemize}

\newpage

% ---------------- TRAINING ----------------
\section*{Тренировочный блок}

\textbf{Распределение по сложности:} 1--3 --- простые, 4--6 --- средние, 7--9 --- выше среднего, 10 --- сложная.

\begin{enumerate}
    \item Первая труба пропускает на $12$ л/мин меньше, чем вторая. Резервуар объёмом $96$ л вторая труба заполняет на $4$ минуты быстрее. Найдите производительность второй трубы.

    \item Первая труба пропускает на $8$ л/мин меньше, чем вторая. Обе заполняют резервуар объёмом $72$ л. Первая труба тратит на $3$ минуты больше. Найдите производительность второй трубы.

    \item Станок~A производит на $5$ деталей в минуту меньше, чем станок~B. На изготовление партии из $120$ деталей станок~A тратит на $4$ минуты больше. Найдите производительность станка~B.

    \item Диагональ параллелограмма делит его угол на $30^\circ$ и $60^\circ$. Найдите отношение сторон параллелограмма.

    \item В прямоугольном треугольнике расположен квадрат со стороной $x$ так, что отрезки на катетах, прилегающие к квадрату, равны $3$ и $12$. Найдите площадь квадрата.

    \item Найдите расстояние от центра ромба до его стороны, если сторона ромба равна $10$, а острый угол равен $60^\circ$.

    \item Первая труба пропускает на $5$ л/мин меньше, чем вторая. Известно, что на заполнение резервуара объёмом $200$ л первая труба тратит на $10$ минут меньше, чем вторая на заполнение резервуара объёмом $320$ л. Найдите производительность второй трубы.

    \item Диагональ параллелограмма делит его угол на $45^\circ$ и $45^\circ$. Найдите отношение соседних сторон параллелограмма.

    \item В прямоугольном треугольнике размещён квадрат. Отрезки на катетах, прилегающие к квадрату, равны $2\sqrt3$ и $6\sqrt3$. Найдите сторону квадрата и его площадь.

    \item Первая труба пропускает на $6$ л/мин меньше, чем вторая. Если первая заполняет резервуар объёмом $420$ л, а вторая --- резервуар объёмом $830$ л, то время работы у них одинаково. Найдите производительность первой и второй труб.
\end{enumerate}

\newpage

% ---------------- ANSWERS ----------------
\section*{Ответы}

\renewcommand{\arraystretch}{1.35}
\begin{table}[h!]
\centering
\caption{Краткие ответы для самопроверки}
\begin{tabularx}{\linewidth}{>{\centering\arraybackslash}p{1.2cm}X}
\toprule
\textbf{\#} & \textbf{Ответ} \\
\midrule
1 & $24$ л/мин \\
2 & $24$ л/мин \\
3 & $20$ деталей/мин \\
4 & $\dfrac12:\dfrac{\sqrt3}{2}=1:\sqrt3$ \\
5 & $36$ \\
6 & $\dfrac{10\sin 60^\circ}{2}=\dfrac{5\sqrt3}{2}$ \\
7 & $20$ л/мин \\
8 & $1:1$ \\
9 & $x=\sqrt{(2\sqrt3)(6\sqrt3)}=6$, площадь $36$ \\
10 & Первая: $12$ л/мин, вторая: $18$ л/мин \\
\bottomrule
\end{tabularx}
\end{table}

\vspace{0.6cm}

\textbf{Напоминание.} Если ответ получился числом, всегда проверьте:
\begin{itemize}
    \item соответствует ли он смыслу задачи;
    \item положительна ли производительность;
    \item не перепутаны ли сторона и площадь;
    \item не потеряна ли единица измерения.
\end{itemize}

\end{document}